\section{Dictamen y limitaciones}
El modelo PostgreSQL representa las cinco entidades mínimas y sus 17 atributos,
mediante columnas directas o relaciones normalizadas. Las dependencias de
fuente, reloj, contrato, tipo y categoría se separan del hecho observado.
Las catorce relaciones satisfacen 3FN respecto de la cobertura funcional
documentada y los dominios SQL tipados. Las dos tablas auxiliares son parte del
modelo definitivo y se incluyen en el E--R completo.

Los tipos administrativos y categorías demo no representan aparatos acreditados
del inventario municipal. Las fechas, estados, responsables y usuarios
sintéticos permanecen identificados. Los vínculos físicos desconocidos no se
completan por inferencia; PG07 conserva esa distinción. Las coordenadas y
asignaciones territoriales no acreditan vigencia histórica individual.

MongoDB conserva su contrato y sus evidencias nativas: lecturas reales, eventos
por regla académica, alertas de replay local e incidentes anonimizados de 2026.
Las consultas por canal o variable no se presentan como identificación de un
dispositivo físico. No se atribuyen alertas oficiales ni causalidad entre períodos.
El laboratorio separado demuestra MG01 por identidad de sensor ficticio y MG07
mediante tipos administrativos, sin equiparar canal y dispositivo. Tres documentos
con texto libre, atributos heterogéneos y revisiones demuestran almacenamiento y
recuperación de información narrativa sintética. La procedencia limita la
interpretación: no acredita dispositivos físicos ni reportes ciudadanos reales.
El almacenamiento distribuido MongoDB permanece sin demostración. Por ese
requisito abierto no se declara el Sprint 1 académicamente completo.

\section{Conclusiones}
La separación entre almacenamiento relacional y vista contractual permite
normalizar PostgreSQL y conservar la identidad compartida con MongoDB. Las
restricciones interrelacionales se mantienen mediante FK y triggers sin repetir
atributos determinados por catálogos. La verificación combina prueba formal de
DF, ejecución nativa, comparación independiente y reconciliación de identidades.

\begin{samepage}
Los mínimos de volumen se evalúan por motor y por generación, sin sumar replay
ni datos administrativos. La ejecución local no acredita el estado de servidores
compartidos. Hadoop/HDFS, Spark, dashboard completo, análisis distribuidos y la
comparación final entre cuatro motores corresponden a la continuación de
noviembre. Esta exclusión no se extiende al almacenamiento distribuido mencionado
en MongoDB: su demostración y modalidad suficiente permanecen abiertas. Las
vistas materializadas y el particionamiento requieren aclaración de alcance.

\end{samepage}

\begin{thebibliography}{99}
\bibitem{enunciado} Instituto Tecnológico de Costa Rica.
\emph{Enunciado Proyecto SmartCity Analytics, II Semestre 2026}.
PDF facilitado al equipo; entrega general del 25/11/2026.
\bibitem{aire} Ayuntamiento de Madrid.
\emph{Calidad del aire: datos horarios}.
\url{https://datos.madrid.es/dataset/201200-0-calidad-aire-horario}.
\bibitem{meteo} Ayuntamiento de Madrid.
\emph{Datos meteorológicos horarios}.
\url{https://datos.madrid.es/dataset/300352-0-meteorologicos-horarios}.
\bibitem{trafico} Ayuntamiento de Madrid.
\emph{Histórico de datos del tráfico}.
\url{https://datos.madrid.es/dataset/208627-0-transporte-ptomedida-historico}.
\bibitem{accidentes} Ayuntamiento de Madrid, Policía Municipal.
\emph{Accidentes de tráfico con detalle}.
\url{https://datos.madrid.es/dataset/300228-0-accidentes-trafico-detalle}.
Snapshot verificado de 2026; datos personales excluidos de la exportación.
\bibitem{contrato} Repositorio SmartCity.
\emph{Contrato implementado de la generación histórica madrid-draft-v1}.
\archivo{research/datasets/scripts/phase5_postgresql.sql};
\archivo{research/datasets/scripts/unified_synthetic.py};
\archivo{docs/equipo/CONTRATO_UNIFICADO.md}, antecedente de esa generación.
\bibitem{permiso} Repositorio SmartCity.
\emph{Registro de comunicación docente del 10/10/2026:
alcance de los datos administrativos sintéticos}.
\archivo{docs/equipo/EVIDENCIA_AUTORIZACION_DATOS_SINTETICOS.md}.
\bibitem{v2} Repositorio SmartCity.
\emph{Implementación y evidencia de la interfaz Madrid v2}.
\archivo{sprint1/postgres/schema_v2_3nf.sql}, vista
\archivo{observation_contract}; adaptación histórica
\archivo{sprint1/postgres/contract_v2.sql};
\archivo{sprint1/basico/preparar_v2.py};
\archivo{research/datasets/results/expanded-v2-2025-01/execution-d26e40706dcd.json}.
\bibitem{cierre} Repositorio SmartCity.
\emph{Validación histórica de la generación anual madrid-draft-v1}.
\archivo{docs/equipo/CIERRE_VALIDACION_SPRINT1.md}.
\bibitem{alcance} Equipo SmartCity.
\emph{Alcance del Sprint 1: registro de la aclaración docente
del 28/10/2026 comunicada al equipo}.
\archivo{docs/equipo/ALCANCE_CONFIRMADO_SPRINT1.md}.
\bibitem{mongocierre} Equipo SmartCity.
\emph{Cierre técnico nativo MongoDB Madrid v2 y resultados verificables}.
\archivo{sprint1/mongodb/evidencia_final_mongodb.md};
\archivo{sprint1/mongodb/evidencias/20261028T055755-3fea400d/comprobacion_cierre.json};
código de la ejecución: \archivo{804d6a551dc84e22d901b64554c30a712b951f77}.
\bibitem{pg3nf} Equipo SmartCity.
\emph{Modelo PostgreSQL normalizado: dependencias funcionales y validación nativa}.
\archivo{sprint1/postgres/impacto_3nf.md};
\archivo{sprint1/postgres/evidencias/3nf/20261028T071812-55791345.json}.
\bibitem{laboratorio} Equipo SmartCity.
\emph{Demostración académica aislada MongoDB: consultas por sensor, tipos y textos narrativos}.
\archivo{sprint1/mongodb/academico/README.md}; ejecución
\archivo{sprint1/mongodb/evidencias/academico/20261028T161057-eff98453/execution.json}.
\end{thebibliography}
